\chapter{Apprendimento su dati strutturati e grafi}
\label{cap:grafi}

\begin{inbreve}
L'ultimo capitolo estende l'apprendimento ai grafi. Formula il problema delle trasduzioni su grafi. Presenta il
\emph{message passing} delle Deep Graph Networks nelle sue due famiglie, a strati e ricorsiva. Discute i problemi
legati alla profondità (over-smoothing, over-squashing, eterofilia), con alcuni strumenti di analisi. Si chiude con
i kernel per strutture.
\end{inbreve}

\begin{attenzione}[Notazione]
In questo capitolo $\x_v$ è l'etichetta (il vettore delle feature) del nodo $v$ e $\h_v^{(l)}$ il suo stato dopo $l$
passi di message passing, con $l=1,\dots,L$. $A$ è la matrice di adiacenza e $\mathcal N(v)$ l'insieme dei vicini di
$v$. Le formule con una matrice $H$ hanno una riga per nodo.
\end{attenzione}

\section{Perché i grafi}

I dati hanno spesso \textbf{relazioni}: astrazioni di immagini e comprensione di scene, social network e knowledge
graph, alberi sintattici del linguaggio, reti di trasporto, piccole molecole (\ingl{drug design}) e reti biologiche,
termini logici, reti cerebrali. Il corso ha seguito una progressione nel dominio di input. Si è passati dai
\textbf{vettori} (reti feedforward) alle \textbf{sequenze} (RNN), e ora alle \textbf{strutture} (alberi, grafi).
L'obiettivo qui è ottenere modelli \textbf{profondi ed efficienti} per i dati strutturati. Si combina la
stratificazione (più livelli di astrazione) con l'efficienza degli approcci randomizzati o incrementali.

\subsection{I grafi considerati}

Si considerano \textbf{grafi etichettati}. Ogni \textbf{nodo} $v$ ha un'\textbf{etichetta vettoriale} $\x_v$, ad
esempio $[1,0,1,0{,}7]$. Gli archi possono avere un'etichetta ed essere orientati, e possono esserci
\textbf{cicli}. $A$ è la \termine{matrice di adiacenza}: $A_{uv}=1$ se c'è un arco tra $u$ e $v$, 0 altrimenti.

\subsection{Il problema della rappresentazione}

Non esisteva un modo \textbf{sistematico}, valido per qualunque task, per estrarre feature da esempi strutturati. È
un caso di representation learning (\rif{sec:representation}) esteso ai domini strutturati.
\begin{itemize}
  \item Le rappresentazioni basate su \textbf{feature} scelte a mano sono incomplete o dipendenti dal task (ad esempio
        gli indici topologici in chimica).
  \item Le rappresentazioni di dimensione fissa (ad esempio la matrice di adiacenza come vettore) hanno tre
        problemi. Soffrono di sovra-dimensionamento o incompletezza: spreco con il padding, perdita di
        informazione. Hanno il problema dell'\textbf{allineamento} tra grafi diversi: quale nodo di un grafo
        corrisponde a quale dell'altro? Dipendono infine dall'ordine dei nodi, e questo rende difficile
        generalizzare.
\end{itemize}

\begin{idea}[Portare i modelli ai dati]
Trattare la natura propria dei dati di input è la chiave per applicare con successo il ML. Invece di \textbf{portare
i dati ai modelli}, si \textbf{portano i modelli ai dati}. Trasformare grafi in vettori o alberi in sequenze porta
infatti problemi di allineamento e perdita di informazione.
\end{idea}

\begin{esempio}[Molecole]
Nel QSPR/QSAR si vuole correlare la struttura chimica delle molecole con le loro proprietà:
$\text{Proprietà} = T(\text{Struttura})$. La proprietà può essere un valore numerico (regressione) o ``tossico
sì/no'' (classificazione). Le molecole non sono vettori: sono rappresentate più naturalmente da grafi di
dimensione variabile, e si vuole predire direttamente dalla struttura.
\end{esempio}

\section{Deep Graph Networks}

\subsection{Trasduzioni su grafi}

Si vuole apprendere una \textbf{trasduzione} $T$ da un dominio di grafi (di dimensione variabile) a uno spazio
discreto o continuo, a partire da esempi $(\text{grafo}_i,\text{target}_i)$. Un encoder $T_{\mathit{enc}}$ calcola un
\textbf{embedding} $\h_v$ di ogni nodo (lo spazio latente). Da qui
\begin{itemize}
  \item un'uscita \textbf{per ogni nodo} dà i compiti \textbf{a livello di nodo}: trasduzioni struttura-struttura
        isomorfe, come la classificazione dei nodi;
  \item una funzione di \textbf{readout} aggrega gli embedding dei nodi in un embedding del grafo $\h_g$. Dà i
        compiti \textbf{a livello di grafo}: regressione o classificazione di grafi;
  \item esistono anche trasduzioni non isomorfe e compiti \textbf{sugli archi}, ad esempio la \ingl{link
        prediction}.
\end{itemize}

\subsection{Il message passing}

Il funzionamento di una \termine{Deep Graph Network} (DGN) si descrive con il \termine{message passing}: i nodi si
scambiano messaggi per diventare consapevoli della struttura circostante (\figrif{fig:message-passing}).
\begin{enumerate}
  \item Ogni nodo calcola un \textbf{messaggio}, funzione della propria etichetta, dei vicini e degli archi, e lo
        invia ai vicini.
  \item I nodi \textbf{aggregano} i messaggi ricevuti con una funzione \textbf{invariante per permutazione}: l'ordine
        di arrivo non conta.
  \item Ogni nodo \textbf{aggiorna} il proprio stato combinando i messaggi con i parametri liberi $W$.
\end{enumerate}
\textbf{Iterando} il message passing l'informazione si \textbf{diffonde} da ogni nodo agli altri. Infine si emette
un'uscita per ogni nodo, oppure per l'intero grafo. Per il grafo si usa un'aggregazione invariante per
permutazione delle rappresentazioni dei nodi, il \termine{global pooling} (ad esempio somma o media).

\begin{figure}[ht]
\centering
\begin{tikzpicture}[x=1cm, y=1cm, nd/.style={circle, draw, fill=grigiochiaro, inner sep=0pt, minimum size=6mm, font=\small}]
\node[nd, fill=rosso!25] (v) at (0,0) {$v$};
\foreach \a/\n in {90/1,200/2,330/3} \node[nd] (u\n) at (\a:1.6) {$u_\n$};
\node[nd] (w1) at (1.6,2.4) {}; \node[nd] (w2) at (-2.4,1.3) {};
\draw (u1) -- (w1); \draw (u2) -- (w2); \draw (u1) -- (w2);
\foreach \n in {1,2,3} \draw[->, very thick, rosso] (u\n) -- (v);
\node[font=\footnotesize, align=left, anchor=west] at (2.3,0.3)
  {1.\ ogni vicino $u\in\mathcal N(v)$ invia un messaggio\\2.\ $v$ aggrega i messaggi (somma, media: invariante\\\phantom{2.\ }per permutazione)\\3.\ $v$ aggiorna il proprio stato $\h_v$ con i pesi $W$};
\end{tikzpicture}
\caption{Il message passing su un nodo $v$: messaggi dai vicini, aggregazione, aggiornamento dello stato.}
\label{fig:message-passing}
\end{figure}

\subsection{Il legame con le CNN}

Sia le CNN sia le reti convoluzionali per grafi visitano ogni posizione dei dati con gli stessi pesi. In una
\textbf{griglia} (un'immagine) però i vicini di un pixel sono \textbf{ordinati e in numero fisso}. In un
\textbf{grafo} sono \textbf{non ordinati e in numero variabile}. Per questo sui grafi si limita il campo recettivo ai
vicini, si condividono i pesi in modo \textbf{invariante per permutazione}, e si estende il vicinato con la
\textbf{stratificazione}.

\subsection{La formula generale}

\begin{equation}\label{eq:dgn}
\boxed{\;\h_v^{(l)} = \mathrm{AGG}_{W^{(l)}}\Big(\x_v,\ \h_v^{(l-1)},\ \{\h_u^{(l-1)} : u\in\mathcal N(v)\}\Big),\qquad l = 1,\dots,L\;}
\end{equation}
dove $\mathrm{AGG}$ aggrega i messaggi dei vicini con un operatore invariante per permutazione (ad esempio una
somma), e $W^{(l)}$ sono i parametri liberi. $\mathcal N$ è dato dalla matrice di adiacenza. Applicare il calcolo a
ogni nodo è una \textbf{visita parallela e non ordinata} del grafo. Alcune varianti, scritte con la matrice
$H^{(l)}$ che ha per riga $v$ lo stato $\h_v^{(l)}$:
\begin{itemize}
  \item \textbf{GCN}: $H^{(l)} = \mathrm{ReLU}\big(\hat A\,H^{(l-1)}W^{(l)}\big)$, con $\hat A$ matrice di adiacenza
        normalizzata per grado;
  \item \textbf{GraphConv}: $H^{(l)} = f\big(H^{(l-1)}W_1^{(l)} + A\,H^{(l-1)}W_2^{(l)}\big)$;
  \item \textbf{GAT}: aggregazione basata sull'attenzione (\rif{sec:transformer}).
\end{itemize}

\begin{idea}[Le DGN sono intrinsecamente profonde]
Il processo non è locale. Iterando sugli strati, la rappresentazione di un nodo include progressivamente
l'informazione di contesto di nodi sempre più lontani (\figrif{fig:contesto}). La \textbf{profondità serve a far
crescere il contesto}.
\end{idea}

\begin{figure}[ht]
\centering
\begin{tikzpicture}[x=0.68cm, y=0.68cm, nd/.style={circle, draw, inner sep=0pt, minimum size=4mm}]
\foreach \s/\r in {0/0,1/1,2/2} {
  \begin{scope}[xshift=\s*5.4cm]
    \coordinate (P0) at (0,0); \coordinate (P1) at (1.2,0.5); \coordinate (P2) at (-1.1,0.7); \coordinate (P3) at (0.2,-1.1);
    \coordinate (P4) at (2.4,0.1); \coordinate (P5) at (-2.2,0.1); \coordinate (P6) at (1.4,-1.4); \coordinate (P7) at (3.3,0.9);
    \foreach \a/\b in {0/1,0/2,0/3,1/4,2/5,3/6,4/7,4/6} \draw[grigio] (P\a) -- (P\b);
    \foreach \k in {0,...,7} \node[nd] at (P\k) {};
    \ifnum\r>0 \foreach \k in {1,2,3} \node[nd, fill=rosso!35] at (P\k) {}; \fi
    \ifnum\r>1 \foreach \k in {4,5,6} \node[nd, fill=rosso!15] at (P\k) {}; \fi
    \node[nd, fill=rosso] at (P0) {};
    \node[font=\footnotesize] at (0.5,-2.2) {strato \the\numexpr\r+1\relax: raggio \r};
  \end{scope}
}
\end{tikzpicture}
\caption{Crescita del contesto di un nodo (rosso) con gli strati di message passing: al primo strato lo stato dipende
dal nodo stesso, poi dai vicini diretti, poi dai vicini dei vicini.}
\label{fig:contesto}
\end{figure}

\subsection{Approccio a strati: GCN e NN4G}

Nelle reti convoluzionali per grafi (GCN, NN4G) \textbf{ogni iterazione del message passing è un nuovo strato}, con
pesi propri. Il contesto di ogni nodo si \textbf{compone} attraverso gli strati: il campo recettivo si estende di
strato in strato. Addestrando il modello si impara come rappresentare il grafo. Le GCN si addestrano end-to-end; NN4G
strato per strato, in modo incrementale. \textbf{NN4G} calcola uno stato scalare per ogni strato e ogni nodo (con
$x_{v,0}=1$ per il bias):
\[
\begin{gathered}
h_v^{(1)} = f\Big(\sum_{j=0}^{n}\bar w_{1j}\,x_{v,j}\Big),\\
h_v^{(l)} = f\left(\underbrace{\sum_{j=0}^{n}\bar w_{lj}\,x_{v,j}}_{\text{etichetta di }v} +
\underbrace{\sum_{i=1}^{l-1}\hat w_{li}\sum_{u\in\mathcal N(v)}h_u^{(i)}}_{\text{vicini, strati precedenti}}\right),\quad l = 2,\dots,L.
\end{gathered}
\]
Non è ricorsiva: ogni strato usa gli stati degli strati \textbf{precedenti}. Le unità si aggiungono una alla volta,
in stile Cascade Correlation (\rif{sec:unita-nascoste}). Il numero di strati è quindi automatico, e non c'è
vanishing gradient attraverso gli strati. Le uscite dei vari strati si combinano per produrre il risultato sul
grafo.

\subsection{Approccio ricorsivo: GNN e GraphESN}

Negli approcci \textbf{ricorsivi} la diffusione è simile, ma $l$ conta le \textbf{iterazioni sullo stesso strato},
con pesi condivisi. È un \textbf{sistema dinamico}. Si richiede la \textbf{convergenza a un punto fisso}, usato come
embedding del grafo. Per questo la dinamica del message passing deve essere \textbf{contrattiva} (\rif{sec:esn}).
Nella \textbf{GNN} l'embedding è addestrato. Nella \textbf{GraphESN} (estensione delle Echo State Network) la funzione
di message passing \textbf{non è addestrata}. È un reservoir casuale di cui si controllano i parametri dinamici:
\[
\h_v^{(l)} = \tanh\Big(W_{\mathit{in}}\,\x_v + \sum_{u\in\mathcal N(v)}\hat W\,\h_u^{(l-1)}\Big),
\]
con $W_{\mathit{in}}$ e $\hat W$ uguali a ogni iterazione. La formula si itera fino a convergenza. L'uscita per un
grafo $g$ si ottiene con un global pooling $\chi$ e un \textbf{readout lineare}, addestrato con pseudoinversa o
ridge regression: $o(g) = W_{\mathit{out}}\,\chi(\h(g))$.

\section{Problemi aperti}

\subsection{Efficienza}

L'addestramento end-to-end con backpropagation è costoso su grafi molto grandi. La GraphESN per la classificazione
di nodi ottiene accuratezze pari o superiori a quelle di molti modelli addestrati end-to-end. Lo fa con
un'accelerazione fino a 10 volte, e scala oltre i 100\,000 nodi. NN4G+ costruisce la rete
\textbf{incrementalmente}, addestrando un'unità alla volta. Uno strato si espande finché l'accuratezza di
validazione migliora, e si aggiungono strati finché migliora. Si ottengono architetture adatte al task in minuti
invece che ore.

\subsection{I problemi della profondità}

\begin{itemize}
  \item \termine{Over-smoothing}: l'accuratezza tipicamente cala oltre 4--5 strati. Il message passing agisce come
        una \textbf{diffusione}, e le rappresentazioni dei nodi negli strati profondi \textbf{collassano},
        diventando tutte simili.
  \item \termine{Over-squashing}: un \textbf{collo di bottiglia}. Il campo recettivo cresce
        \textbf{esponenzialmente} con la profondità, ma tutta quell'informazione va compressa in un embedding di
        dimensione fissa per nodo. Si misura la \textbf{sensibilità} della rappresentazione di un nodo all'input di
        nodi lontani con lo Jacobiano $\partial\h_v^{(L)}/\partial\x_u$ (\rif{sec:tecniche-deep}). Questa
        sensibilità decresce esponenzialmente con il numero di strati.
  \item Per l'addestramento end-to-end c'è la difficoltà di \textbf{retropropagare il gradiente} attraverso molti
        strati.
\end{itemize}
Gli effetti sono due. Il primo è l'\termine{under-reaching}: non si ottiene il campo recettivo necessario per
apprendere interazioni tra nodi lontani. Il secondo è la difficoltà con l'\termine{eterofilia}. Nei grafi
\textbf{omofili} i nodi vicini condividono per lo più la stessa classe. In quelli \textbf{eterofili} i nodi della
stessa classe sono in genere lontani, e predizioni basate sui vicini immediati possono essere fuorvianti
(\figrif{fig:omofilia}).

\begin{figure}[ht]
\centering
\begin{tikzpicture}[x=0.85cm, y=0.85cm, nd/.style={circle, draw, inner sep=0pt, minimum size=4.2mm}]
\foreach \s/\cols in {0/{blu,blu,blu,blu,rosso,rosso,rosso,rosso}, 1/{blu,rosso,blu,rosso,rosso,blu,rosso,blu}} {
  \begin{scope}[xshift=\s*7cm]
    \coordinate (Q1) at (0,1); \coordinate (Q2) at (1,1.8); \coordinate (Q3) at (1.2,0.4); \coordinate (Q4) at (2.2,1.2);
    \coordinate (Q5) at (3.4,1.4); \coordinate (Q6) at (4.3,2.0); \coordinate (Q7) at (4.4,0.6); \coordinate (Q8) at (5.3,1.3);
    \foreach \a/\b in {1/2,1/3,2/3,2/4,3/4,4/5,5/6,5/7,6/8,7/8,6/7} \draw[grigio] (Q\a) -- (Q\b);
    \foreach \c [count=\k] in \cols \node[nd, fill=\c!60] at (Q\k) {};
  \end{scope}
}
\node[font=\footnotesize] at (2.65,-0.5) {omofilia alta};
\node[font=\footnotesize] at (9.65,-0.5) {omofilia bassa (eterofilia)};
\end{tikzpicture}
\caption{Grafi ad alta e bassa omofilia (il colore è la classe): i secondi sono difficili per il message passing
standard.}
\label{fig:omofilia}
\end{figure}

\subsection{Strumenti di analisi}

Modelli come NN4G e GraphESN permettono di \textbf{separare il bias del message passing dai problemi
dell'addestramento end-to-end}. NN4G è incrementale e non propaga il gradiente attraverso gli strati. GraphESN non
addestra affatto l'embedding.

Nella GraphESN si controlla direttamente la \textbf{norma spettrale} $\norm{W}$ dei pesi. È il massimo fattore di
cui la matrice può allungare un vettore. Coincide con la \textbf{costante di Lipschitz} della funzione di message
passing, cioè il massimo rapporto tra la variazione dell'uscita e la variazione dell'input. Vale
\[
\left\lVert\frac{\partial\h_v^{(L)}}{\partial\x_u}\right\rVert \le \prod_{l=1}^{L}\norm{W^{(l)}}\cdot\big(A^L\big)_{u,v}.
\]
È un limite \textbf{superiore}. Se ogni $\norm{W^{(l)}}$ è minore di 1, il prodotto decresce esponenzialmente con
$L$: la sensibilità rispetto all'input dei nodi lontani è \textbf{costretta} a svanire (over-squashing). Imporre
$\norm W>1$ toglie questo vincolo: il limite non obbliga più la sensibilità a svanire esponenzialmente. Non
garantisce però che non svanisca per altri motivi. C'è anche una tensione con la contrattività chiesta sopra
all'approccio ricorsivo. La contrattività ($\norm W<1$) garantisce la convergenza a un punto fisso, ma è proprio la
condizione che costringe la sensibilità a svanire. Controllare direttamente $\norm W$ permette di studiare questo
compromesso.

Un secondo strumento è l'\textbf{analisi spettrale}. La ``frequenza'' di un segnale su un grafo (etichette,
embedding) è la sua variabilità tra nodi vicini. Si misura con l'\termine{energia di Dirichlet}
\[
E(\h) = \sum_v\sum_{u\in\mathcal N(v)}\norm{\h_v-\h_u}^2 .
\]
L'energia dice quanto il segnale varia sul grafo: un segnale liscio ha bassa energia. Si usa anche il
\termine{Laplaciano} del grafo, $\mathcal L = D-A$, con $D$ matrice diagonale dei gradi (il numero di vicini di
ogni nodo). Dalla decomposizione di $\mathcal L$ nei suoi autovettori si ottiene la trasformata di Fourier su grafo.
Dice come il segnale si distribuisce sulle frequenze. I grafi eterofili hanno segnali ad \textbf{alta frequenza}, e
i task eterofili richiedono di \textbf{preservare le alte frequenze}. Su un benchmark eterofilo (una rete di pagine
Wikipedia), NN4G e GraphESN preservano meglio l'energia attraverso gli strati. Mantengono l'accuratezza
all'aumentare dei passi di message passing, mentre quella di GCN, GAT e GraphConv cala.

\subsection{Estensioni}

\begin{itemize}
  \item Grafi \textbf{dinamici}, con relazioni o valori che evolvono nel tempo, trattati ad esempio con reservoir
        multistrato.
  \item \textbf{Generazione} di grafi: apprendere la distribuzione $p(G)$ per campionarne di nuovi, ad esempio
        molecole.
  \item \textbf{Spiegabilità}: individuare le sotto-strutture che determinano una predizione.
  \item \textbf{Network quantification}: predire la prevalenza delle classi tra i nodi non etichettati.
  \item Applicazioni in \textbf{bioinformatica}.
\end{itemize}

\section{Kernel per strutture}
\label{sec:kernel-strutture}

I metodi kernel (\rif{sec:metodi-kernel}) si estendono facilmente a sequenze, alberi e grafi. Basta definire un
prodotto scalare $k(s,s') = \langle\phi(s),\phi(s')\rangle$, cioè una \textbf{similarità} tra oggetti $s,s'$ di
qualsiasi tipo, con un embedding implicito in uno spazio euclideo.
\begin{itemize}
  \item \textbf{Kernel marginalizzati}: la similarità tra due grafi si basa sui cammini di etichette comuni,
        ottenuti con random walk. Il kernel è il prodotto scalare dei vettori di conteggio, mediato su tutti i
        cammini. Il costo è tipicamente quadratico nella dimensione del grafo.
  \item \textbf{Kernel di convoluzione}: si decompongono gli oggetti in \textbf{sotto-strutture}, e si definisce il
        kernel combinando kernel tra le sotto-strutture.
\end{itemize}

\begin{attenzione}[Limiti dei kernel per strutture]
\begin{itemize}
  \item \textbf{Efficienza}: la scomposizione in sotto-strutture rende il costo critico.
  \item \textbf{Adattività}: la similarità è \textbf{fissata prima dell'apprendimento}, quindi la codifica non è
        appresa. Il vantaggio è che si può usare la conoscenza a priori sul problema.
  \item Per i domini strutturati non si conosce l'\textbf{universalità}: la scelta del kernel è critica. Calcolare
        un kernel \textbf{completo} per grafi è difficile almeno quanto il problema dell'\textbf{isomorfismo di
        grafi}.
\end{itemize}
\end{attenzione}

\begin{idea}[Conclusione]
Se i dati hanno relazioni (seriali, gerarchiche, reti), non conviene usare una rappresentazione piatta. Il deep
learning su grafi si può fare anche \textbf{in modo efficiente}, senza un addestramento completamente end-to-end.
La profondità nelle DGN è necessaria, perché fa crescere il contesto. Causa però problemi (over-smoothing,
over-squashing, difficoltà di addestramento) che sono oggetto di ricerca.
\end{idea}

\begin{puntichiave}
\begin{itemize}
  \item Trasduzioni su grafi: livello di nodo, di grafo, di arco.
  \item Message passing: messaggio, aggregazione invariante per permutazione, aggiornamento;
        $\h_v^{(l)} = \mathrm{AGG}_{W^{(l)}}(\x_v,\h_v^{(l-1)},\{\h_u^{(l-1)}\}_{u\in\mathcal N(v)})$.
  \item A strati (GCN, NN4G: ogni iterazione è uno strato) o ricorsivo (GNN, GraphESN: iterazioni contrattive fino al
        punto fisso).
  \item Over-smoothing, over-squashing, under-reaching, eterofilia.
  \item Kernel per strutture: marginalizzati, di convoluzione; similarità fissata a priori.
\end{itemize}
\end{puntichiave}

\begin{domande}
\begin{enumerate}
  \item Perché trattare i dati strutturati direttamente invece di trasformarli in vettori?
  \item Quali trasduzioni si possono definire sui grafi?
  \item Descrivere il message passing e la formula generale delle DGN.
  \item Differenza tra approcci a strati (GCN, NN4G) e ricorsivi (GNN, GraphESN).
  \item Cosa sono over-smoothing e over-squashing? Cos'è l'eterofilia?
  \item Pro e contro dei kernel per strutture.
\end{enumerate}
\end{domande}
